\documentclass[fleqn]{beamer}

\mode<presentation> {
\usetheme{Madrid}
\usecolortheme{dolphin}
}

\usepackage[utf8]{inputenc}
\usepackage[brazilian]{babel}
\usepackage{amsmath}
\usepackage{amssymb}
\usepackage{booktabs}

\title[Aula 2 -- Problema da Dieta]{O Problema da Dieta}
\subtitle{O desafio da nutricionista}
\author[Santi]{Prof. \'Everton Santi}
\institute[UFRN/ECT]{ECT2524 -- Introdução à Otimização \\ Escola de Ciências e Tecnologia -- UFRN}
\date{Aula 2}

\begin{document}

\begin{frame}
	\titlepage
\end{frame}

\begin{frame}
	\frametitle{Sumário}
	\tableofcontents
\end{frame}

%------------------------------------------------------------
\section{Recapitulando}
%------------------------------------------------------------

\begin{frame}
	\frametitle{Na aula passada: Problema da Mochila}
	\begin{itemize}
		\item Escolher itens \textbf{binários} (usa ou não usa) que
			\textbf{maximizam} um benefício sob um limite de capacidade;
		\item Formulamos, escrevemos em AMPL (\texttt{.mod}, \texttt{.dat},
			\texttt{.run}) e resolvemos com HiGHS;
		\item Hoje: um problema com a mesma estrutura de modelagem, mas
			decisões \textbf{contínuas} e vários limites em vez de um só.
	\end{itemize}
\end{frame}

%------------------------------------------------------------
\section{O Desafio}
%------------------------------------------------------------

\begin{frame}
	\frametitle{Estudo de caso: o refeitório da fábrica}
	\begin{itemize}
		\item Uma fábrica serve \textbf{milhares de refeições por dia}
			no seu refeitório;
		\item Isso representa um gasto \textbf{mensal alto} com
			alimentação;
		\item A empresa quer \textbf{reduzir o custo} das refeições --
			mas \textbf{sem prejudicar a saúde} dos trabalhadores;
		\item A nutricionista responsável topa ajudar, mas o problema
			tem \textbf{muitas combinações possíveis} de alimentos e
			quantidades para avaliar ``na mão'';
		\item Ela recorre ao setor de planejamento (vocês) para montar
			um modelo de otimização que sugira o cardápio de menor
			custo.
	\end{itemize}
\end{frame}

\begin{frame}
	\frametitle{O que é o Problema da Dieta?}
	\begin{itemize}
		\item Existe um conjunto de \textbf{alimentos disponíveis}, cada
			um com um \textbf{custo} e uma composição nutricional
			conhecida (calorias, proteínas, vitaminas, etc.);
		\item É preciso decidir \textbf{quais alimentos usar e em que
			quantidade} para compor a refeição/cardápio;
		\item A escolha precisa respeitar \textbf{exigências
			nutricionais}: quantidades mínimas de certos nutrientes e um
			limite máximo de calorias, por exemplo;
		\item Entre todas as combinações que atendem essas exigências,
			queremos a de \textbf{menor custo};
		\item É um problema clássico de otimização, estudado desde a
			década de 1940 (planejamento de dietas militares) e ainda
			usado hoje em nutrição, agropecuária e logística.
	\end{itemize}
\end{frame}

\begin{frame}
	\frametitle{Exigências nutricionais do estudo de caso}
	\begin{itemize}
		\item A nutricionista fixou, por dia, quantidades
			\textbf{mínimas} de:
		\begin{itemize}
			\item proteínas, cálcio, sódio, ferro e vitaminas;
		\end{itemize}
		\item E um limite \textbf{máximo} de calorias;
		\item Cada alimento contribui de forma diferente para cada um
			desses nutrientes -- por isso a combinação importa, não só
			a quantidade total de comida.
	\end{itemize}
\end{frame}

%------------------------------------------------------------
\section{Formulação}
%------------------------------------------------------------

\begin{frame}
	\frametitle{Problema da Dieta -- semelhança com a mochila}
	\begin{itemize}
		\item Mochila: \textbf{maximiza} benefício sob um limite (peso);
		\item Dieta: \textbf{minimiza} custo sob \textbf{vários} limites
			(um por nutriente), com mínimo e máximo;
		\item Mochila: variáveis \textbf{binárias};
		\item Dieta: variáveis \textbf{contínuas} não-negativas;
		\item Mesma estrutura de modelagem -- muda o que se decide.
	\end{itemize}
\end{frame}

%------------------------------------------------------------
\section{Tarefas da Aula 2}
%------------------------------------------------------------

\begin{frame}
	\frametitle{Duas tarefas nesta aula}
	\begin{enumerate}
		\item \textbf{Tarefa 2a} -- a dieta de \textbf{um único dia},
			sem se preocupar com estoque de alimentos;
		\item \textbf{Tarefa 2b} -- a mesma dieta, mas considerando
			\textbf{estoque} de alimentos, \textbf{total de refeições}
			por dia e um \textbf{horizonte de 5 dias} (uma semana
			útil).
	\end{enumerate}
	\vspace{0.2cm}
	A Tarefa 2b usa a formulação e os dados da Tarefa 2a como ponto de
	partida -- resolvam-nas nessa ordem.
\end{frame}

\begin{frame}
	\frametitle{Tarefa 2a -- Problema da Dieta (um dia, sem estoque)}
	\begin{enumerate}
		\item Pesquisem, por conta própria, dados nutricionais reais de
			alimentos (tabela TACO, rótulos, bases nutricionais, etc.);
		\item Montem \textbf{seu próprio} \texttt{dieta.dat}, com pelo
			menos 10 alimentos e as propriedades nutricionais pedidas;
		\item Escrevam a formulação compacta do problema (parâmetros,
			variáveis, objetivo, restrições);
		\item Implementem em AMPL (\texttt{.mod}, \texttt{.dat},
			\texttt{.run}) e resolvam com HiGHS;
	\end{enumerate}
	\vspace{0.2cm}
	Detalhes completos na página da aula, no site do curso.
\end{frame}

\begin{frame}
	\frametitle{Tarefa 2b -- Dieta com estoque (horizonte de 5 dias)}
	\begin{enumerate}
		\item Reaproveitem os alimentos e dados da Tarefa 2a;
		\item Definam um \textbf{estoque inicial} de cada alimento e o
			\textbf{total de refeições} a servir em cada um dos 5 dias;
		\item Estendam a formulação: um índice a mais para o
			\textbf{dia}, restrições de \textbf{balanço de estoque}
			entre os dias, e exigências nutricionais multiplicadas
			pelo total de refeições de cada dia;
		\item Implementem em AMPL (\texttt{dieta\_estoque.mod/.dat/.run})
			e resolvam com HiGHS, para os 5 dias;
	\end{enumerate}
	\vspace{0.2cm}
	Detalhes completos na página da aula, no site do curso.
\end{frame}

\begin{frame}
	\frametitle{Perguntas para reflexão}
	\begin{enumerate}
		\item Quais as limitações da solução obtida?
		\item O que poderia ser melhorado no modelo?
		\item Como pensar o cardápio de um refeitório real (ex.: RU do
			campus), ao longo de uma semana, mês ou ano?
		\item Este modelo poderia ser aplicado na sua casa ou trabalho?
	\end{enumerate}
\end{frame}

\begin{frame}
	\frametitle{Referências}
	\begin{itemize}
		\item AMPL: \url{https://ampl.com};
	\end{itemize}
\end{frame}

\end{document}
